
\subsubsection{Laplace and infinite $\E[R^2]$}
We will see now that if $q$ is Laplace in \textbf{dimension one}, then $\E[R^2]$ is infinite.
\begin{theorem}
    \label{prop:laplace_infinite_second_moment}
    Assume $d=1$, and let $q(t)=\frac12 e^{-|t|}$,
    i.e., $T\sim \mathrm{Laplace}(0,1)$.
    Then
    \begin{align}
        \E[R_0] = \E[R^2]=\infty\,,
    \end{align}
    for every $\lambda \in (0,1)$.
\end{theorem}
The proof is deferred to \cref{proof:laplace_infinite_second_moment}.
Together with \cref{prop:exact_quadratic_forces_chi2}, this shows that
\begin{align}
    \lim_{\alpha\downarrow0}
    \frac{1}{\alpha^2}
    D_{f_\alpha}(\PP_{X, T_0} || \PP_X \otimes \PP_{T_0})
    =
    \infty\,,
\end{align}
\paragraph{Lesson:}
it is impossible to get exact $\alpha^2$ scaling of the variance term
for a general log-concave $q$ (e.g., Laplace).

However, this is unclear if one can circumvent the Laplace example by either (i) putting further assumptions on $q$ (e.g., subgaussian tails), or (ii) paying a $\polylog(1/\alpha)$ factor.
We first study (ii).

\subsection{Paying log factors through $f_\alpha$ via reduction to tail bounds}
\label{sec:log_factor_reduction_to_tail}
The next result shows that if $q$ is Laplace in dimension one, then we can still get the desired $\tilde{O}(\eta^2)$ bound by paying a logarithmic factor.
This also shows that $f_\alpha$-divergence is crucial and one would lose a lot (literally infinite)
if simply using the $\chi^2$-divergence as a surrogate for $f_\alpha$-divergence.
We prove this through a sequence of lemmas that reduce the problem of bounding $D_{f_\alpha}(\PP_{X, T_0} || \PP_X \otimes \PP_{T_0})$ to bounding the tail probability of $R_0$.
Recall
\begin{align}
    R_0=\frac{P_{X,T_0}(X, T_0)}{P_X(X)P_{T}(T_0)}\,, \qquad
    \text{and}\qquad
    R=\frac{P_{X,T_0}(X, T)}{P_X(X)P_T(T)} \qquad \text{for } T \stackrel{i.i.d}{\sim} p_{T_0}\,,
\end{align}

\begin{lemma}[$f_\alpha$ is qaudratic-then-linear]
    \label{lem:log_factor_fa_quadratic_then_linear}
    For every $\alpha\in(0,1)$ and $P\ll Q$ we have
    \begin{align}
        D_{f_\alpha}(P\|Q)
         & \le
        e^{\alpha}
        \alpha^2
        \left(
        1+\E[R_0\mathbf{1}_{\{R_0\le \frac{1}{\alpha}\}}]
        \right)
        +
        \alpha \PP\left(R_0>\frac{1}{\alpha}\right)\,.
        \label{eq:log_factor_reduction_to_tail}
    \end{align}
\end{lemma}
The proof is deferred to \cref{proof:log_factor_fa_quadratic_then_linear}.

The next result shows that how a tail bound on $R_0$ implies a bound on $D_{f_\alpha}(\PP_{X, T_0} || \PP_X \otimes \PP_{T_0})$.
\begin{lemma}[Tail to divergence control]
    \label{lem:log_factor_tail_to_divergence}
    Suppose that for every $y\ge1$,
    \marginpar{\textit{Note that if we had $\E[R_0] = C <\infty$, then this bound would be trivial.
            But $\E[R_0]$ is the chi-squared divergence, which is infinite.}}
    \begin{align}
        \PP(R_0>y)
        \le
        \frac{C}{y}\,,
        \label{eq:log_factor_R0_tail}
    \end{align}
    for some constant $C>0$.
    Then for every $\alpha\in(0,1)$,
    \begin{align}
        D_{f_\alpha}(P\|Q)
         & \le
        C \alpha^2\log\left(\frac{e}{\alpha}\right).
        \label{eq:log_factor_fa_upper}
    \end{align}
\end{lemma}
The proof is deferred to \cref{proof:log_factor_tail_to_divergence}.

\marginpar{Basically we are almost sure that $R_0$ has a Pareto tail.}
Now, it is clear that if we can show the tail bound in \eqref{eq:log_factor_R0_tail} for $R_0$, then we get the desired $\tilde{O}(\eta^2)$ bound on the variance term.
We first do it for the Laplace distribution, and then we show that the same tail bound holds for a general log-concave distribution, in 1 dimension.

\noindent
First, define
\begin{align}
    X'\triangleq\lambda\sum_{k=1}^\infty (1-\lambda)^{k-1} T_k,
    \qquad
    \text{so that }\qquad
    X=\lambda T_0+(1-\lambda) X'.
\end{align}
Note that $X' \stackrel{d.}{\equiv} X$ and $X' \perp T_0$.
Let $T \stackrel{i.i.d.}{\sim} T_0$; the following identity for $R_0$ plays a crucial role in the analysis:
\begin{align}
    R_0
     & =
    \frac{p_{X|T_0}(X)}{p_X(X)}
    =
    \frac{p_{X|T_0}(\lambda T_0 + (1-\lambda) X')}{p_{X}(X)}
    =
    \frac{1}{(1-\lambda)^d}\cdot \frac{p_X(X')}{p_X(X)}
    \label{eq:R0_identity_1}
    \\
     & =
    \frac{1}{(1-\lambda)^d}\cdot \frac{p_X(X')}{p_X(\lambda T_0+(1-\lambda)X')}
    =
    \frac{1}{(1-\lambda)^d}\cdot \frac{p_X(X)}{p_X(\lambda T+(1-\lambda)X)}\,.
    \label{eq:R0_identity_2}
\end{align}

\subsubsection{Bounds in 1d}
We now show that how the tail bounds on $R_0$ can be established in the 1-dimensional case when $q$ is Laplace, and then we show that the same tail bound holds for a general log-concave distribution in 1 dimension.

\begin{theorem}[Laplace bound]
    \label{theorem:laplace_log_factor}
    Assume $d=1$, and let $q(t)=\frac12 e^{-|t|}$.
    Then there exists a numerical constant $C>0$ such that for every
    $\lambda\in(0,1)$ and every $\alpha\in(0,1]$,
    \begin{align}
        D_{f_\alpha}(\PP_{X,T_0}\|\PP_X\otimes\PP_{T_0})
        \le
        \frac{4}{1-\lambda}
        \alpha^2 \log\left(\frac{e}{\alpha}\right)\,.
        \label{eq:laplace_log_factor_fa_upper}
    \end{align}
    % In particular,
    % \begin{align}
    %     \sup_{f:\K\to[0,1]}
    %     \Lambda_{\mathrm{var},q}(E_\lambda,p_X;f)
    %     \le
    %     \frac{C}{1-\lambda}
    %     \frac{\eta^2}{\lambda^2}
    %     \log\left(\frac{e\lambda}{\eta}\right)
    %     +
    %     \frac{\eta^2}{2}.
    %     \label{eq:laplace_log_factor_variance_upper}
    % \end{align}
    % Moreover, the logarithm is not an artefact: there exist constants
    % $c_\lambda>0$ and $a_\lambda\in(0,1]$, depending only on $\lambda$,
    % such that for every $a\in(0,a_\lambda]$,
    % \begin{align}
    %     D_{f_a}(\PP_{X,T_0}\|\PP_X\otimes\PP_{T_0})
    %     \ge
    %     c_\lambda a^2 \log\left(\frac1a\right).
    %     \label{eq:laplace_log_factor_fa_lower}
    % \end{align}
\end{theorem}
% The proof is deferred to \cref{proof:laplace_log_factor}.
% \subsubsection{Proof of \cref{theorem:laplace_log_factor}}
% \label{proof:laplace_log_factor}
\begin{proof}
    For a 1-dimensional application of \cref{eq:R0_identity_2}, we have
    \begin{align}
        R_0
         & =
        \frac{1}{1-\lambda}\cdot \frac{p_X(X)}{p_X(\lambda T+(1-\lambda)X)}\,.
        \label{eq:laplace_log_factor_R0_formula}
    \end{align}
    Next, we will show that $p_X$ is log-Lipschitz with Lipschitz constant $1/\lambda$.
    Let
    \begin{align}
        q_\lambda(u)=\frac{1}{2\lambda}e^{-|u|/\lambda}
    \end{align}
    be the density of $\lambda T_0$.
    Let $p_z$ be the density of $Z = (1-\lambda)X'$,
    then since $p=q_\lambda*p_z$, we have the lower bound
    \marginpar{
        \textit{There is an upper bound too:}
        \begin{align}
            \frac{p(x+h)}{p(x)}
            \le
            e^{|h|/\lambda}\,.
        \end{align}
    }
    \begin{align}
        p(x+h)
         & =
        \int_\R q_\lambda(x+h-z)p_z(z)\,dz
        \\
         & \ge
        e^{-|h|/\lambda}
        \int_\R q_\lambda(x-z)p_z(z)\,dz
        =
        e^{-|h|/\lambda}p(x)\,,
        \label{eq:laplace_log_factor_log_lipschitz}
    \end{align}
    which implies that $p$ is $\frac{1}{\lambda}$-log-Lipschitz.
    On the other hand, using the same self-identical decomposition,
    \begin{align}
        p(\lambda t+(1-\lambda)z)
         & =
        \frac{1}{1-\lambda}
        \int_\R
        q(s)
        p\left(
        z+\frac{\lambda(t-s)}{1-\lambda}
        \right)
        ds
        \\
         & \ge
        \frac{p(z)}{1-\lambda}
        \int_\R
        q(s)
        \exp\left(
        -\frac{|t-s|}{1-\lambda}
        \right)
        ds,
        \label{eq:laplace_log_factor_denominator_lower}
    \end{align}
    where we used \eqref{eq:laplace_log_factor_log_lipschitz}.
    Thus, we have
    \begin{align}
        R_0
        =
        \frac{1}{1-\lambda}\cdot \frac{p_X(X)}{p_X(\lambda T+(1-\lambda)X)}
         & \le
        \left(
        \int_\R
        q(s)
        \exp\left(
        -\frac{|T-s|}{1-\lambda}
        \right)
        \d s
        \right)^{-1}
        \label{eq:laplace_log_factor_R0_integral_lower}
        \\
         & =
        \E\left[
            \exp\left(
            -\frac{|T - T'|}{1-\lambda}
            \right)
            \middle | T
            \right]^{-1}\,.
        \label{eq:laplace_log_factor_R0_integral_upper}
    \end{align}
    Since $q(s)=\frac12e^{-|s|}$, for every $t\in\R$,
    \begin{align}
        \E\left[
            \exp\left(
            -\frac{|T - T'|}{1-\lambda}
            \right)
            \middle | T
            \right]
         & =
        \int_\R
        q(s)
        \exp\left(
        -\frac{|T-s|}{1-\lambda}
        \right)
        \d s
        \\
         & =
        2(1-\lambda)
        \int_\R
        q(s)
        \ell_{1-\lambda}(T-s)
        \d s
        \\
         & =
        2(1-\lambda) q * \ell_{1-\lambda}(T)
        \\
         & =
        2(1-\lambda) p_{U}(T)
        \label{eq:laplace_log_factor_laplace_convolution_lower}
    \end{align}
    where $\ell_{1-\lambda}(v)$ is the density of $(1-\lambda)L$ where $L\sim\mathrm{Lap}(0,1)$, and $p_U$ is the density of $U = T + (1-\lambda)L$.
    So we have
    \begin{align}
        R_0
         & \le
        \frac{1}{2(1-\lambda) p_U(T)}
        \quad
        \Rightarrow
        \quad
        \PP(R_0>y)
        \le
        \PP\left(p_U(T)<\frac{1}{2(1-\lambda)y}\right).
    \end{align}

    \marginpar{All the steps so far hold for a general \textbf{globally log-Lipschitz} density $q$ through rescaling.}
    \paragraph{Laplace specific computation.}
    When $t\ge0$, integration gives
    \begin{align}
        \int_\R
        q(s)
        \exp\left(
        -\frac{|t-s|}{1-\lambda}
        \right)
        \d s
         & =
        \frac{1}{2}
        \int_\R
        \exp\left(
        -|s|
        -\frac{|t-s|}{1-\lambda}
        \right)
        \d s    \\
         & \geq
        \frac{1}{2}
        \int_t^\infty
        \exp\left(
        -s
        -\frac{s-t}{1-\lambda}
        \right)
        \d s    \\
         & =
        \frac{1}{2}
        \int_t^{\infty}
        \exp\left(
        -\frac{(2- \lambda) s - t}{1-\lambda}
        \right)
        \d s
        \\
         & =
        \frac{1}{2}
        \int_{0}^\infty
        \exp\left(
        -\frac{(2 - \lambda) (t+s) - t}{1-\lambda}
        \right)
        \d s
        \\
         & =
        \frac{1}{2}
        \int_{0}^\infty
        \exp\left(
        -t - \frac{(2-\lambda)s}{1-\lambda}
        \right)
        \d s
        \\
         & =
        \frac{1}{2}
        \exp(-t)
        \int_{0}^\infty
        \exp\left(
        - \frac{2-\lambda}{1-\lambda} s
        \right)
        \d s
        \\
         & =
        \frac{1-\lambda}{4-2\lambda}
        \exp(-t)\,,
    \end{align}
    and the case $t<0$ is symmetric.
    Combining the last display with \cref{eq:laplace_log_factor_R0_integral_lower} gives
    \begin{align}
        R_0
        \le
        \left(
        \int_\R q(s) \exp\left(- \frac{|T-s|}{1-\lambda}\right) \d s
        \right)^{-1}
        \le
        \frac{4-2\lambda}{1-\lambda}e^{|T|}
        \le
        \frac{4}{1-\lambda}e^{|T|}
        \,.
        \label{eq:laplace_log_factor_R0_exp_upper}
    \end{align}
    Finally, $|T|\sim\mathrm{Exp}(1)$, so
    \begin{align}
        \PP\left(R_0>y\right)
         & \le
        \PP\left(e^{|T|}>\frac{1-\lambda}{4}y\right)\le \frac{4}{(1-\lambda)y}.
    \end{align}
    The result follows from \cref{lem:log_factor_tail_to_divergence}.
\end{proof}

\begin{theorem}[Log-concave kernels pay at most one logarithm]
    \label{theorem:log_concave_log_factor}
    Assume $d=1$, and let $q$ be a log-concave probability density on an interval of $\R$.
    Then there exists a numerical constant $C>0$ such that for every
    $\lambda\in(0,1)$ and every $a\in(0,1]$,
    \begin{align}
        D_{f_a}(\PP_{X,T_0}\|\PP_X\otimes\PP_{T_0})
        \le
        \frac{C}{1-\lambda}
        a^2 \log\left(\frac{e}{a}\right).
        \label{eq:log_concave_log_factor_fa_upper}
    \end{align}
    In particular, for every $0<\eta\le \lambda$,
    \begin{align}
        \sup_{f:\K\to[0,1]}
        \Lambda_{\mathrm{var},q}(E_\lambda,p_X;f)
        \le
        \frac{C}{1-\lambda}
        \frac{\eta^2}{\lambda^2}
        \log\left(\frac{e\lambda}{\eta}\right)
        +
        \frac{\eta^2}{2}.
        \label{eq:log_concave_log_factor_variance_upper}
    \end{align}
\end{theorem}

The proof is deferred to \cref{proof:log_concave_log_factor}.
\subsubsection{Proof of \cref{theorem:log_concave_log_factor}}
\label{proof:log_concave_log_factor}
\begin{proof}
    We prove that
    \begin{align}
        \PP(R_0>y)
        \le
        \frac{C}{(1-\lambda)y},
        \qquad y>0.
        \label{eq:log_concave_log_factor_R0_tail}
    \end{align}
    Let $p=p_X$ be the density of $X$.
    Since $q$ is log-concave and log-concavity is preserved by affine transformations,
    convolution, and weak limits, $p$ is also log-concave.

    Define
    \begin{align}
        X'=\lambda\sum_{k=1}^\infty (1-\lambda)^{k-1} T_k,
        \qquad
        \text{so that }\qquad
        X=\lambda T_0+(1-\lambda) X'.
    \end{align}
    Note that $X' \stackrel{d.}{\equiv} X$ and $X' \perp T_0$,
    so we have
    \begin{align}
        R_0
        =
        \frac{p_{X|T_0}(X)}{p_X(X)}
         & =
        \frac{p_{X|T_0}(\lambda T_0 + (1-\lambda) X')}{p_{X}(X)}
        =
        \frac{1}{1-\lambda}\cdot \frac{p_X(X')}{p_X(X)}\,,\label{eq:log_concave_log_factor_R0_formula}
    \end{align}
    Let $\|p\|_\infty = \sup_{x\in\R} p(x)$.
    From \eqref{eq:log_concave_log_factor_R0_formula}, we get the pointwise bound
    \marginpar{\textit{This inequality is essentially very crude, and potentially should be improved to get the right rate in higher dimensions.}}
    \begin{align}
        R_0
        \le
        \frac{1}{1-\lambda}\frac{\|p\|_\infty}{p(X)}.
        \label{eq:log_concave_log_factor_R0_density_upper}
    \end{align}
    Thus it remains to control how often a one-dimensional log-concave density is much smaller than its maximum.
    We use the following consequence of log-concavity.
    \begin{lemma}[1D Log-concave density concentration]
        Let $Y$ have log-concave density $h$, and let $\|h\|_\infty=h(m)$ for a mode $m$.
        Then, for every $\epsilon > 0$,
        \begin{align}
            \PP\left(h(Y)\le \epsilon \|h\|_\infty\right)
            \le
            2\epsilon.
            \label{eq:log_concave_density_small_ball}
        \end{align}
    \end{lemma}
    \begin{proof}
        First, the result is trivial if $\epsilon\ge1$, so we assume $\epsilon<1$.
        To see this, first consider the right tail. If $x\ge m$, then log-concavity implies that
        for every $u\ge x$,
        \begin{align}
            \frac{h(u)}{h(u-(x-m))}
            \le
            \frac{h(x)}{h(m)}.
        \end{align}
        Hence
        \begin{align}
            \PP(Y\ge x)
            =
            \int_x^\infty h(u)\,du
            \le
            \frac{h(x)}{h(m)}
            \int_x^\infty h(u-(x-m))\,du
            \le
            \frac{h(x)}{\|h\|_\infty}.
            \label{eq:log_concave_right_tail_density_bound}
        \end{align}
        The same argument on the left gives
        \begin{align}
            \PP(Y\le x)
            \le
            \frac{h(x)}{\|h\|_\infty},
            \qquad x\le m.
            \label{eq:log_concave_left_tail_density_bound}
        \end{align}
        Since the sublevel set
        \begin{align}
            \{x:h(x)\le \epsilon\|h\|_\infty\}
        \end{align}
        is contained in the union of a left tail and a right tail, each having probability at most
        $\epsilon$, \eqref{eq:log_concave_density_small_ball} follows.
    \end{proof}

    Applying \eqref{eq:log_concave_density_small_ball} with $h=p$ and
    \begin{align}
        \epsilon=\frac{1}{(1-\lambda)t},
    \end{align}
    we obtain
    \begin{align}
        \PP(R_0>t)
         & \le
        \PP\left(
        p(X)<\frac{\|p\|_\infty}{(1-\lambda)t}
        \right) \le
        \min\left\{
        1,
        \frac{2}{(1-\lambda)t}
        \right\}.
    \end{align}
    This proves \eqref{eq:log_concave_log_factor_R0_tail}.
\end{proof}

\subsection{Log-concave measures in higher dimensions}
\begin{lemma}[Truncated exponential moment of a log-concave density]
    \label{lem:log_concave_truncated_exp_potential}
    Let $h$ be a log-concave probability density on $\Rd$. Write
    \begin{align}
        h(x)=\|h\|_\infty e^{-V(x)},
    \end{align}
    where $V:\Rd\to[0,\infty]$ is convex and $\inf_x V(x)=0$.
    Let $Y\sim h$. Then there exists a numerical constant $C>0$ such that
    for every $L\ge0$,
    \begin{align}
        \E\left[
            \min\left\{\frac{\|h\|_\infty}{h(Y)},e^L\right\}
            \right]
        =
        \E\left[
            \min\left\{e^{V(Y)},e^L\right\}
            \right]
        \le
        C
        \sum_{j=0}^d \frac{L^j}{j!}.
        \label{eq:log_concave_truncated_exp_potential}
    \end{align}
\end{lemma}

\begin{proof}
    Let
    \begin{align}
        S_d(L)=\sum_{j=0}^d\frac{L^j}{j!}.
    \end{align}
    The result is trivial for $L\le d$, since
    \begin{align}
        \E\left[
            \min\left\{e^{V(Y)},e^L\right\}
            \right]
        \le e^L
        \le 2 S_d(L),
    \end{align}
    where the last inequality is the standard Poisson-tail estimate
    $e^L\le 2\sum_{j=0}^d L^j/j!$ for $L\le d$.

    We now consider $L>d$.
    After translating, assume that the minimum of $V$ is attained at $0$.
    For $s\ge0$, define
    \begin{align}
        K_s=\{x\in\Rd:V(x)\le s\},
        \qquad
        F(s)=\mathrm{vol}(K_s).
    \end{align}
    By convexity of $V$, for every $0<s\le t$,
    \begin{align}
        \frac{s}{t}K_t\subseteq K_s.
    \end{align}
    Hence
    \begin{align}
        \frac{F(s)}{s^d}
        \ge
        \frac{F(t)}{t^d},
        \qquad 0<s\le t.
        \label{eq:sublevel_volume_ratio_monotone}
    \end{align}
    First, by integration by parts, we have
    \begin{align}
        1
        \geq
        \int_{K_L} h(x)\d x
         & =
        \|h\|_\infty
        \int_{V(x)\le L} e^{-V(x)}\d x \\
         & =
        \|h\|_\infty
        \int_0^L e^{-s}\d F(s)         \\
         & =
        \|h\|_\infty
        \left(
        e^{-L}F(L)+\int_0^L e^{-s}F(s)\d s
        \right)                        \\
         & \geq
        \|h\|_\infty
        \left(
        e^{-L}F(L)+\int_0^L e^{-s} \left(\frac{s}{L}\right)^d  F(L)\d s
        \right)                        \\
         & =
        \|h\|_\infty
        \frac{F(L)}{L^d}
        \left(
        L^de^{-L}+\int_0^L e^{-s} s^d\d s
        \right)                        \\
         & =
        \|h\|_\infty
        \frac{F(L)}{L^d}
        \left(
        d \int_0^L e^{-s} s^{d-1}\d s
        \right)\,,
    \end{align}
    and since $L>d$, the last integral is at least a numerical constant times
    $d!$. Thus
    \begin{align}
        \|h\|_\infty F(L)
        \le
        C\frac{L^d}{d!}.
        \label{eq:log_concave_sublevel_volume_factorial}
    \end{align}

    Next, for $s\ge L$, \eqref{eq:sublevel_volume_ratio_monotone} gives
    \begin{align}
        F(s)\le \left(\frac{s}{L}\right)^d F(L).
    \end{align}
    Therefore, again by integration by parts,
    \begin{align}
        \PP(V(Y)>L)
         & =
        \|h\|_\infty
        \int_L^\infty e^{-s}\d F(s)
        \le
        \|h\|_\infty
        \frac{F(L)}{L^d}
        d\int_L^\infty e^{-s}s^{d-1}\d s.
    \end{align}
    Combining this with \eqref{eq:log_concave_sublevel_volume_factorial},
    and using the exact formula for the Gamma tail,
    \begin{align}
        d\int_L^\infty e^{-s}s^{d-1}\d s
        =
        d!e^{-L}\sum_{j=0}^{d-1}\frac{L^j}{j!},
    \end{align}
    gives
    \begin{align}
        e^L\PP(V(Y)>L)
        \le
        C\sum_{j=0}^{d-1}\frac{L^j}{j!}.
        \label{eq:log_concave_potential_tail_factorial}
    \end{align}
    Finally,
    \begin{align}
        \E\left[
            \min\left\{e^{V(Y)},e^L\right\}
            \right]
         & =
        \E\left[
            e^{V(Y)}\mathbf{1}_{\{V(Y)\le L\}}
            \right]
        +
        e^L\PP(V(Y)>L)
        \\
         & =
        \|h\|_\infty F(L)
        +
        e^L\PP(V(Y)>L)
        \\
         & \le
        C
        \left(
        \frac{L^d}{d!}
        +
        \sum_{j=0}^{d-1}\frac{L^j}{j!}
        \right)
        =
        C
        \sum_{j=0}^d\frac{L^j}{j!}.
    \end{align}
    This proves the result.
\end{proof}


\begin{theorem}[Log-concave kernels pay \(d\) logarithms with factorial dependence]
    \label{theorem:log_concave_d_log_factor_factorial}
    Assume that $q$ is a log-concave probability density on $\Rd$.
    Then there exists a numerical constant $C>0$ such that for every
    $\lambda\in(0,1)$ and every $a\in(0,1]$,
    \begin{align}
        D_{f_a}(\PP_{X,T_0}\|\PP_X\otimes\PP_{T_0})
        \le
        \frac{C}{(1-\lambda)^d}
        a^2
        \sum_{j=0}^d
        \frac{1}{j!}
        \log^j\left(\frac{e}{a}\right).
        \label{eq:log_concave_d_log_factor_factorial_fa_upper}
    \end{align}
    In particular, for every $0<\eta\le \lambda$,
    \begin{align}
        \sup_{f:\K\to[0,1]}
        \Lambda_{\mathrm{var},q}(E_\lambda,p_X;f)
        \le
        \frac{C}{(1-\lambda)^d}
        \frac{\eta^2}{\lambda^2}
        \sum_{j=0}^d
        \frac{1}{j!}
        \log^j\left(\frac{e\lambda}{\eta}\right)
        +
        \frac{\eta^2}{2}.
        \label{eq:log_concave_d_log_factor_factorial_variance_upper}
    \end{align}
\end{theorem}

\begin{proof}
    Suppose that $d>1$, then
    \begin{align}
        B_\lambda\triangleq\frac{1}{(1-\lambda)^d}\,,
        \qquad
        \text{and when }\lambda=\frac{1}{d}, \qquad e < B_\lambda < 2e\,.
    \end{align}
    % \paragraph{Case 1: $aB_\lambda\ge1$.}
    % We have
    % \begin{align}
    %     \E\left[
    %         R_0\mathbf{1}_{\{R_0\le 1/a\}}
    %         \right]
    %     \le
    %     \frac1a,
    %     \qquad
    %     \PP\left(R_0>\frac1a\right)\le1.
    % \end{align}
    % Hence
    % \begin{align}
    %     D_{f_a}(\PP_{X,T_0}\|\PP_X\otimes\PP_{T_0})
    %     \le
    %     C a.
    % \end{align}
    % Since $aB_\lambda\ge1$, we have $a\le a^2B_\lambda$, so
    % \begin{align}
    %     D_{f_a}(\PP_{X,T_0}\|\PP_X\otimes\PP_{T_0})
    %     \le
    %     C a^2B_\lambda
    %     \le
    %     C a^2B_\lambda
    %     \sum_{j=0}^d
    %     \frac{1}{j!}
    %     \log^j\left(\frac{e}{a}\right).
    %     \label{eq:log_concave_d_factorial_large_a}
    % \end{align}

    % \paragraph{Case 2: $aB_\lambda<1$.}
    % Suppose that $aB_\lambda<1$, the other case is trivial since $a\le a^2B_\lambda$.
    Recall the self-identical decomposition
    \begin{align}
        X'=\lambda\sum_{k=1}^\infty (1-\lambda)^{k-1}T_k,
        \qquad
        X=\lambda T_0+(1-\lambda)X'\,.
    \end{align}
    Since $q$ is log-concave, $p$ is log-concave.
    Write
    \begin{align}
        p(x)=\|p\|_\infty e^{-V(x)}
    \end{align}
    with $V$ convex and $\inf_x V(x)=0$. From \eqref{eq:R0_identity_1},
    \begin{align}
        R_0
        =
        B_\lambda
        \frac{p(X')}{p(X)}
        \le
        B_\lambda
        \frac{\|p\|_\infty}{p(X)}
        =
        B_\lambda e^{V(X)}.
        \label{eq:log_concave_d_factorial_R0_upper}
    \end{align}

    We use the same quadratic-then-linear reduction as in the one-dimensional
    proof:
    \begin{align}
        D_{f_a}(\PP_{X,T_0}\|\PP_X\otimes\PP_{T_0})
        \le
        e^a
        a^2
        \left(
        1+
        \E\left[
            R_0\mathbf{1}_{\{R_0\le 1/a\}}
            \right]
        \right)
        +
        a\PP\left(R_0>\frac1a\right).
        \label{eq:log_concave_d_factorial_reduction}
    \end{align}
    Suppose that $aB_\lambda<1$ since the other case is trivial since $a\le a^2B_\lambda$, and define
    \begin{align}
        L=\log\left(\frac{1}{aB_\lambda}\right)>0.
    \end{align}
    From \eqref{eq:log_concave_d_factorial_R0_upper},
    \begin{align}
        \E\left[
            R_0\mathbf{1}_{\{R_0\le 1/a\}}
            \right]
         & \le
        \E\left[
            \min\left\{R_0,\frac1a\right\}
            \right]
        \le
        B_\lambda
        \E\left[
            \min\left\{e^{V(X)},e^L\right\}
            \right].
    \end{align}
    Applying \cref{lem:log_concave_truncated_exp_potential} with $h=p$ and
    $Y=X$ gives
    \begin{align}
        \E\left[
            R_0\mathbf{1}_{\{R_0\le 1/a\}}
            \right]
        \le
        C B_\lambda
        \sum_{j=0}^d \frac{L^j}{j!}.
        \label{eq:log_concave_d_factorial_truncated_mean}
    \end{align}
    Similarly,
    \begin{align}
        \PP\left(R_0>\frac1a\right)
         & \le
        \PP\left(B_\lambda e^{V(X)}>\frac1a\right)
        =
        \PP(V(X)>L).
    \end{align}
    Since
    \begin{align}
        e^L\PP(V(X)>L)
        \le
        \E\left[
            \min\left\{e^{V(X)},e^L\right\}
            \right],
    \end{align}
    another application of \cref{lem:log_concave_truncated_exp_potential}
    yields
    \begin{align}
        a\PP\left(R_0>\frac1a\right)
         & \le
        ae^{-L}
        C
        \sum_{j=0}^d \frac{L^j}{j!}
        =
        C a^2B_\lambda
        \sum_{j=0}^d \frac{L^j}{j!}.
        \label{eq:log_concave_d_factorial_tail_term}
    \end{align}
    Plugging \eqref{eq:log_concave_d_factorial_truncated_mean} and
    \eqref{eq:log_concave_d_factorial_tail_term} into
    \eqref{eq:log_concave_d_factorial_reduction}, and using
    $B_\lambda\ge1$, $a\le1$, and $L\le \log(e/a)$, gives
    \begin{align}
        D_{f_a}(\PP_{X,T_0}\|\PP_X\otimes\PP_{T_0})
         & \le
        C a^2B_\lambda
        \sum_{j=0}^d
        \frac{1}{j!}
        \log^j\left(\frac{e}{a}\right)\,.
    \end{align}
\end{proof}

The natural question is whether the $d$-logarithm factor is real or an artifact of the proof.
\begin{theorem}[Product Laplace forces the Gamma factor]
    \label{thm:product_laplace_gamma_lower}
    Fix $\lambda\in(0,1)$. For each $d\ge1$, let
    \[
        \ell(t)=\frac12 e^{-|t|},\qquad t\in\R,
        \qquad
        q_d(t)=\prod_{i=1}^d \ell(t_i),
        \qquad t=(t_1,\ldots,t_d)\in\R^d .
    \]
    Let $T_0,T_1,\ldots$ be i.i.d. samples from $q_d$, and define
    \[
        X=\lambda\sum_{k=0}^{\infty}(1-\lambda)^kT_k .
    \]
    Let
    \[
        P=P_{X,T_0},
        \qquad
        Q=P_X\otimes P_{T_0},
        \qquad
        R=\frac{dP}{dQ}.
    \]
    For $a>0$, define
    \[
        f_a(u)=e^{-a(u-1)}+a(u-1)-1,\qquad u\ge0,
    \]
    and
    \[
        D_{f_a}(P\|Q)=\E_Q[f_a(R)].
    \]
    Then there exist constants
    \[
        A_\lambda>0,\qquad
        \tau_\lambda<\infty,\qquad
        \beta_\lambda>0,
    \]
    depending only on $\lambda$, with $A_\lambda e^{\tau_\lambda}\ge2$, such that the following holds.
    With the universal constant $c_0=e^{-1}/8$, for every $d\ge1$ and every
    \[
        0<a\le (A_\lambda e^{\tau_\lambda})^{-d},
    \]
    we have
    \[
        D_{f_a}(P\|Q)
        \ge
        c_0
        \left(\frac{\beta_\lambda A_\lambda}{2}\right)^d
        a^2
        \sum_{j=0}^d
        \frac{L_a^j}{j!},
    \]
    where
    \[
        L_a
        =
        \log\left(
        \frac{1}{a(A_\lambda e^{\tau_\lambda})^d}
        \right).
    \]
    Consequently, for each fixed $\lambda\in(0,1)$ and $d\ge1$, there exist constants
    $c_{\lambda,d}>0$ and $a_{\lambda,d}\in(0,1)$ such that for all
    $0<a\le a_{\lambda,d}$,
    \[
        D_{f_a}(P\|Q)
        \ge
        c_{\lambda,d}
        a^2
        \sum_{j=0}^d
        \frac{\log^j(1/a)}{j!}.
    \]
    In particular,
    \[
        D_{f_a}(P\|Q)
        =
        \Omega_{\lambda,d}
        \left(
        a^2\frac{\log^d(1/a)}{d!}
        \right)
        \qquad\text{as }a\downarrow0.
    \]
\end{theorem}

\begin{proof}
    The proof isolates the source of the $d$ logarithms.  In one dimension,
    if the first Laplace sample is large and positive while the remaining
    contribution is bounded, then the likelihood ratio grows like $e^t$.
    In $d$ dimensions this effect tensorizes, so the likelihood ratio grows
    like $\exp(\sum_i t_i)$.  The Laplace density contributes
    $\exp(-\sum_i t_i)$, and truncating at height $1/a$ leaves exactly a
    Gamma-tail polynomial.

    \paragraph{Step 1: the one-dimensional bad event.}
    Let $L_0,L_1,\ldots$ be i.i.d. $\mathrm{Laplace}(0,1)$ random variables, and define
    \[
        X_\lambda
        =
        \lambda\sum_{k=0}^{\infty}(1-\lambda)^kL_k,
        \qquad
        X_\lambda'
        =
        \lambda\sum_{k=1}^{\infty}(1-\lambda)^{k-1}L_k .
    \]
    Let $p_\lambda$ be the density of $X_\lambda$. Then
    $X_\lambda'\stackrel{d}{=}X_\lambda$, $X_\lambda'\perp L_0$, and
    \[
        X_\lambda=\lambda L_0+(1-\lambda)X_\lambda'.
    \]
    Therefore, at a point of the form
    \[
        X_\lambda=\lambda t+(1-\lambda)z,
        \qquad
        L_0=t,
        \qquad
        X_\lambda'=z,
    \]
    the one-dimensional likelihood ratio is
    \[
        r_\lambda(t,z)
        =
        \frac{1}{1-\lambda}
        \frac{p_\lambda(z)}
        {p_\lambda(\lambda t+(1-\lambda)z)}.
        \label{eq:product_laplace_one_dim_ratio}
    \]

    We now prove that this ratio is exponentially large when $t$ is large
    and $z$ is bounded.  Write
    \[
        Z_\lambda
        =
        \lambda\sum_{k=1}^{\infty}(1-\lambda)^kL_k,
        \qquad
        X_\lambda=\lambda L_0+Z_\lambda.
    \]
    If
    \[
        q_\lambda(u)=\frac{1}{2\lambda}e^{-|u|/\lambda}
    \]
    is the density of $\lambda L_0$, then
    \[
        p_\lambda=q_\lambda*p_{Z_\lambda}.
    \]
    Define
    \[
        C_\lambda
        =
        \E\left[e^{|Z_\lambda|/\lambda}\right].
    \]
    This constant is finite, since
    \[
        \frac{|Z_\lambda|}{\lambda}
        \le
        \sum_{k=1}^{\infty}(1-\lambda)^k|L_k|,
    \]
    and hence, using $\E[e^{\alpha |L_0|}]=(1-\alpha)^{-1}$ for
    $0\le\alpha<1$,
    \[
        C_\lambda
        \le
        \prod_{k=1}^{\infty}
        \E\left[e^{(1-\lambda)^k|L_k|}\right]
        =
        \prod_{k=1}^{\infty}
        \frac{1}{1-(1-\lambda)^k}
        <\infty.
    \]
    Consequently, for every $x\ge0$,
    \[
        p_\lambda(x)
        =
        \int_\R
        \frac{1}{2\lambda}e^{-|x-z|/\lambda}
        p_{Z_\lambda}(z)\,dz
        \le
        \frac{C_\lambda}{2\lambda}e^{-x/\lambda}.
        \label{eq:product_laplace_density_tail}
    \]
    Since $p_\lambda$ is positive and continuous, also define
    \[
        m_\lambda
        =
        \inf_{|z|\le1}p_\lambda(z)
        >
        0.
    \]

    Now fix $|z|\le1$ and $t\ge(1-\lambda)/\lambda$. Then
    $\lambda t+(1-\lambda)z\ge0$, so \eqref{eq:product_laplace_density_tail} gives
    \[
        p_\lambda(\lambda t+(1-\lambda)z)
        \le
        \frac{C_\lambda}{2\lambda}
        \exp\left(
        -t-\frac{1-\lambda}{\lambda}z
        \right)
        \le
        \frac{C_\lambda}{2\lambda}
        \exp\left(
        -t+\frac{1-\lambda}{\lambda}
        \right).
    \]
    Substituting this into \eqref{eq:product_laplace_one_dim_ratio} yields
    \[
        r_\lambda(t,z)
        \ge
        \frac{2\lambda m_\lambda}{(1-\lambda)C_\lambda}
        \exp\left(-\frac{1-\lambda}{\lambda}\right)e^t .
    \]
    Set
    \[
        A_\lambda
        =
        \frac{2\lambda m_\lambda}{(1-\lambda)C_\lambda}
        \exp\left(-\frac{1-\lambda}{\lambda}\right).
    \]
    Choose $\tau_\lambda<\infty$ large enough that
    \[
        \tau_\lambda\ge\frac{1-\lambda}{\lambda}
        \qquad\text{and}\qquad
        A_\lambda e^{\tau_\lambda}\ge2.
    \]
    Then
    \[
        r_\lambda(t,z)\ge A_\lambda e^t,
        \qquad
        t\ge\tau_\lambda,\quad |z|\le1.
        \label{eq:product_laplace_one_dim_ratio_lower}
    \]
    Finally define
    \[
        \beta_\lambda
        =
        \PP(|X_\lambda'|\le1)
        =
        \int_{-1}^1 p_\lambda(z)\,dz
        >
        0.
    \]

    \paragraph{Step 2: tensorize the bad event.}
    Return to dimension $d$. Define
    \[
        X'
        =
        \lambda\sum_{k=1}^{\infty}(1-\lambda)^{k-1}T_k.
    \]
    Then $X'\stackrel{d}{=}X$, $X'\perp T_0$, and
    \[
        X=\lambda T_0+(1-\lambda)X'.
    \]
    Under $P=P_{X,T_0}$, write
    \[
        R_0=R(X,T_0)=\frac{dP}{dQ}(X,T_0).
    \]
    Since $q_d$ is a product distribution and the construction of $X$ is coordinatewise,
    the density of $X$ factorizes as
    \[
        p_X(x)=\prod_{i=1}^d p_\lambda(x_i).
    \]
    Therefore,
    \[
        R_0
        =
        \prod_{i=1}^d
        r_\lambda(T_{0,i},X_i').
        \label{eq:product_laplace_ratio_tensorization}
    \]

    Consider the event
    \[
        \mathcal G
        =
        \left\{
        |X_i'|\le1 \text{ and } T_{0,i}\ge\tau_\lambda
        \text{ for every } i=1,\ldots,d
        \right\}.
    \]
    On this event, \eqref{eq:product_laplace_one_dim_ratio_lower} and
    \eqref{eq:product_laplace_ratio_tensorization} imply
    \[
        R_0
        \ge
        A_\lambda^d
        \exp\left(\sum_{i=1}^d T_{0,i}\right).
        \label{eq:product_laplace_R0_large}
    \]
    Moreover, since $A_\lambda e^{\tau_\lambda}\ge2$, the same event implies
    \[
        R_0
        \ge
        (A_\lambda e^{\tau_\lambda})^d
        \ge2.
        \label{eq:product_laplace_R0_at_least_two}
    \]

    \paragraph{Step 3: turn the bad event into an $f_a$-divergence lower bound.}
    We first record an elementary lower bound on the generator.  For every
    $a\in(0,1]$ and every $u\ge2$,
    \[
        \frac{f_a(u)}{u}
        \ge
        c_0 a^2\min\left\{u,\frac1a\right\},
        \qquad
        c_0=\frac{e^{-1}}{8}.
        \label{eq:product_laplace_generator_lower}
    \]
    To prove this, let $s=a(u-1)$. If $s\le1$, then
    \[
        e^{-s}+s-1\ge \frac{e^{-1}}{2}s^2,
    \]
    and since $u\ge2$,
    \[
        \frac{(u-1)^2}{u}\ge \frac{u}{4}.
    \]
    Hence
    \[
        \frac{f_a(u)}{u}
        \ge
        \frac{e^{-1}}{8}a^2u
        \ge
        \frac{e^{-1}}{8}a^2\min\left\{u,\frac1a\right\}.
    \]
    If $s>1$, then it is enough to prove a lower bound of order $a$.
    When $1<s\le2$,
    \[
        f_a(u)=e^{-s}+s-1\ge e^{-1},
        \qquad
        u=\frac{s+a}{a}\le \frac{3}{a},
    \]
    so $f_a(u)/u\ge e^{-1}a/3$. When $s\ge2$,
    \[
        f_a(u)=e^{-s}+s-1\ge s-1\ge \frac{s}{2},
        \qquad
        u=\frac{s+a}{a}\le \frac{3s}{2a},
    \]
    so $f_a(u)/u\ge a/3$. This proves
    \eqref{eq:product_laplace_generator_lower}.

    Since $R=dP/dQ$, the change of measure identity gives
    \[
        D_{f_a}(P\|Q)
        =
        \E_Q[f_a(R)]
        =
        \E_P\left[\frac{f_a(R_0)}{R_0}\right].
    \]
    By the assumption $a\le(A_\lambda e^{\tau_\lambda})^{-d}$ and the fact
    $A_\lambda e^{\tau_\lambda}\ge2$, we have $a\le1$. Therefore, using
    \eqref{eq:product_laplace_generator_lower} on the event $\mathcal G$,
    where $R_0\ge2$, gives
    \[
        D_{f_a}(P\|Q)
        \ge
        c_0a^2
        \E_P\left[
            \min\left\{R_0,\frac1a\right\}
            \mathbf{1}_{\mathcal G}
            \right].
    \]
    Using \eqref{eq:product_laplace_R0_large}, and using the independence of
    $X'$ and $T_0$, we get
    \[
        \begin{aligned}
            D_{f_a}(P\|Q)
             & \ge
            c_0a^2\beta_\lambda^d
            \int_{[\tau_\lambda,\infty)^d}
            \min\left\{
            A_\lambda^d
            \exp\left(\sum_{i=1}^d t_i\right),
            \frac1a
            \right\}
            \prod_{i=1}^d
            \left(\frac12 e^{-t_i}\right)
            dt_i .
        \end{aligned}
        \label{eq:product_laplace_integral_lower}
    \]

    It remains only to evaluate the integral in
    \eqref{eq:product_laplace_integral_lower}. Set
    \[
        L_a
        =
        \log\left(
        \frac{1}{a(A_\lambda e^{\tau_\lambda})^d}
        \right).
    \]
    The assumption on $a$ gives $L_a\ge0$.  Put
    \[
        t_i=\tau_\lambda+u_i,
        \qquad
        u_i\ge0,
        \qquad
        S=\sum_{i=1}^d u_i.
    \]
    Then the integral in \eqref{eq:product_laplace_integral_lower} equals
    \[
        \begin{aligned}
            I_a
             & =
            2^{-d}
            \int_{\R_+^d}
            \min\left\{
            A_\lambda^d e^{d\tau_\lambda}e^S,
            \frac1a
            \right\}
            e^{-d\tau_\lambda}e^{-S}
            du
            \\
             & =
            2^{-d}
            \left[
                A_\lambda^d
                \operatorname{vol}\{u\in\R_+^d:S\le L_a\}
                +
                \frac{e^{-d\tau_\lambda}}{a}
                \int_{\{u\in\R_+^d:S>L_a\}} e^{-S}\,du
                \right].
        \end{aligned}
    \]
    The first term is a simplex volume:
    \[
        \operatorname{vol}\{u\in\R_+^d:S\le L_a\}
        =
        \frac{L_a^d}{d!}.
    \]
    For the second term, using the Gamma tail identity,
    \[
        \int_{\{u\in\R_+^d:S>L_a\}} e^{-S}\,du
        =
        \int_{L_a}^{\infty}
        e^{-s}\frac{s^{d-1}}{(d-1)!}\,ds
        =
        e^{-L_a}\sum_{j=0}^{d-1}\frac{L_a^j}{j!}.
    \]
    Since
    \[
        e^{-L_a}=a(A_\lambda e^{\tau_\lambda})^d,
    \]
    we obtain
    \[
        \frac{e^{-d\tau_\lambda}}{a}
        e^{-L_a}
        =
        A_\lambda^d.
    \]
    Hence
    \[
        I_a
        =
        2^{-d}A_\lambda^d
        \sum_{j=0}^d
        \frac{L_a^j}{j!}.
    \]
    Plugging this value of $I_a$ into
    \eqref{eq:product_laplace_integral_lower} gives
    \[
        D_{f_a}(P\|Q)
        \ge
        c_0
        \left(\frac{\beta_\lambda A_\lambda}{2}\right)^d
        a^2
        \sum_{j=0}^d
        \frac{L_a^j}{j!},
    \]
    which proves the main lower bound.

    \paragraph{Step 4: derive the asymptotic form.}
    For fixed $\lambda$ and $d$,
    \[
        L_a
        =
        \log(1/a)
        -
        d\log(A_\lambda e^{\tau_\lambda}).
    \]
    Therefore, if
    \[
        0<a\le
        a_{\lambda,d}
        \triangleq
        (A_\lambda e^{\tau_\lambda})^{-2d},
    \]
    then
    \[
        L_a\ge \frac12\log(1/a).
    \]
    Consequently,
    \[
        \sum_{j=0}^d\frac{L_a^j}{j!}
        \ge
        2^{-d}
        \sum_{j=0}^d
        \frac{\log^j(1/a)}{j!}.
    \]
    Absorbing the fixed factor
    \[
        c_0
        \left(\frac{\beta_\lambda A_\lambda}{2}\right)^d
        2^{-d}
    \]
    into $c_{\lambda,d}$ gives
    \[
        D_{f_a}(P\|Q)
        \ge
        c_{\lambda,d}
        a^2
        \sum_{j=0}^d
        \frac{\log^j(1/a)}{j!}
    \]
    for all $0<a\le a_{\lambda,d}$. The leading term in the sum is
    $\log^d(1/a)/d!$, which gives
    \[
        D_{f_a}(P\|Q)
        =
        \Omega_{\lambda,d}
        \left(
        a^2\frac{\log^d(1/a)}{d!}
        \right)
        \qquad\text{as }a\downarrow0.
    \]
\end{proof}

\newpage
\section{Ensuring log-concavity of p}
There are various steps in the previous section, where one can get a bound on the variance term, which is effectively $\frac{1}{\eta^2}D_{f_\alpha}(\PP_{X,T_0}\|\PP_X\otimes\PP_{T_0})$ given \textit{some kind of information} about the law of $T$, which is $q$.
For instance in \cref{theorem:log_concave_d_log_factor_factorial}, we showed a bound that needed $q$ to be log-concave.
Or one could get a bound by departing from the step \cref{eq:laplace_log_factor_R0_integral_upper} in the proof of the 1d laplace, which will easily generalizes to higher dimensions.

The outstanding question is
\begin{center}
    \textit{What can be ensured about $q$ that allows us to get a bound on the variance term?}
\end{center}
Note that $q$ is the exponential weights distribution on the set of historic estimated losses $\hf_t$, namely
\begin{align}
    q(t) = \frac{e^{-\eta \hf_t}}{\int_\K e^{-\eta \hf_s} ds}\,,
\end{align}
where $\hf_s = \cE_{X_s}^{\ell_s} = \cE(X_s, \ell_s(X_s))$ is the estimated loss in round $s$,
after taking action $X_s$ and observing the loss $\ell_s(X_s)$.
So, for instance, if the $\hf_t$ are convex functions of $t$, then $q$ is log-concave.

However, our choice of the estimator $\cE_x^\ell(t) = \ell(x) s_x(t)$, where
\begin{align}
    s_x(t) = 1 + \frac{r(x,t) - 1}{\lambda},
\end{align}
where $r(x,t) = \frac{p(x|t)}{p(x)}$ is the likelihood ratio, and all quantities are defined with respect to the state of the world at time step $t$ (i.e., the distribution of $X_s$ and the conditional distribution of $X_s$ given $T=t$), does not necessarily give us convex $\hf_t$, although, the expectation of $\hf_t$ is the convex function $\tilde{\ell}_t = \E_{X_s}[\cE_{X_s}^{\ell_s}(t)]$.

The first hope is that, even if each $\hf_t$ is not convex, after enough rounds, the exponential weights distribution $q$ which is based on the $\hf_t$ will be \textbf{roughly} log-concave, since it is based on the sum of many $\hf_t$.

The reason for being \textbf{roughly} log-concave is that the proof of the log-concave case only needs the $p$ to be log-concave, and $p$ is the convolution of scaled versions of $q$, so even if $q$ is not exactly log-concave, but is close enough to a log-concave distribution, then $p$ might still be log-concave, and we can still get the desired bound on the variance term.

Research this question and come back with an answer!



\subsubsection{Proof of \cref{prop:laplace_infinite_second_moment}}
\label{proof:laplace_infinite_second_moment}
\begin{proof}
    Recall that
    \begin{align}
        R_0
        = \frac{p_{X,T_0}(X, T_0)}{p_X(X) p_{T_0}(T_0)}
        = \frac{p_{X|T_0=T_0}(X)}{p_X(X)}\,.
    \end{align}
    The lower bound idea is to consider a very large value of $x$, say $x=10000$, then $p_X(10000)$ is very small as it is really a tail event.
    However, if we condition on $T_0=10000/\lambda$, then $p_{X|T_0=10000/\lambda}(10000)$ is not a tail event at all, and it is much bigger than $p_X(10000)$, so the ratio is very large.
    \textbf{The main saving grace here is that the probability of $t=10000/\lambda$ is also very small}, but it turns out that it is not small enough to cancel this out.

    \paragraph{Step 1: the bad strip.}
    However, we show that $R_0$ is actually exponentially large when
    \begin{align}
        X \in [\lambda T_0 - \epsilon, \lambda T_0 + \epsilon]\,, \qquad\text{and}\qquad T_0 > 2\epsilon/\lambda\,,
    \end{align}
    and $\epsilon$ is a small constant, say $\epsilon = \frac12$.

    \paragraph{Step 2: $p(x)$ has a $\lambda$-Laplace tail.}
    Define $Z$ as
    \begin{align}
        Z = \lambda\sum_{k=1}^\infty (1-\lambda)^k T_k\,,
        \qquad\text{so that}\qquad
        X = \lambda T_0 + Z\,,
    \end{align}
    where $Z \perp T_0$ and $Z \stackrel{d.}{=}(1-\lambda)X$,
    so its density is given by
    \begin{align}
        p_z(z) = \frac{1}{1-\lambda} p_X\left(\frac{z}{1-\lambda}\right)\,.
    \end{align}
    Let
    \begin{align}
        q_\lambda(u)\triangleq \frac{1}{2\lambda}e^{-|u|/\lambda},
    \end{align}
    which is the density of $\lambda T_0$.
    Since $X=\lambda T_0+Z$, we have $p=q_\lambda * p_z$, that is,
    \begin{align}
        p(x)=\int_{\R} q_\lambda(x-z)p_z(z)\,dz.
    \end{align}
    Therefore, for $x>0$,
    \begin{align}
        p(x)
        =
        \int_{\R}\frac{1}{2\lambda}e^{-|x-z|/\lambda}p_z(z)\,dz
        \le
        \frac{1}{2\lambda}e^{-x/\lambda}\int_{\R}e^{|z|/\lambda}p_z(z)\,dz\,,
    \end{align}
    where we used $|x-z|\ge x-|z|$.
    Set
    \begin{align}
        C_\lambda \triangleq \int_{\R} e^{|z|/\lambda}p_z(z)\,dz = \E[e^{|Z|/\lambda}]\,.
    \end{align}
    We claim that $C_\lambda<\infty$ and we prove it at the end.
    So the tail bound is
    \begin{align}
        p(x) \le \frac{C_\lambda}{2\lambda}e^{-x/\lambda}\,.
        \label{eq:laplace_p_tail_bound}
    \end{align}

    \paragraph{Step 3: lower bound the integrand on an infinite strip.}
    First, note that
    \begin{align}
        \E[R_0]
         & =
        \int_{\R}\int_{\R}
        \left(
        \frac{P_{X,T_0}(x, t)}{P_X(x) P_{T_0}(t)}
        \right)
        P_{X,T_0}(x, t) dx\,dt
        \\
         & =
        \int_{\R}\int_{\R}
        \frac{P_{X|T_0=t}(x)}{P_X(x)}
        P_{X,T_0}(x, t) dx\,dt
        \\
         & =
        \int_{\R}\int_{\R}
        \frac{q(t)}{p(x)} P_{X|T_0=t}(x)^2 dx\,dt
        \\
         & \ge
        \frac{\lambda}{C_\lambda}
        \int_{\R}\int_{x>0}
        \,\frac{
            e^{-|t|}
        }{e^{-x/\lambda}}
        P_{X|T_0=t}(x)^2
        \,dx\,dt
        \\
         & \geq
        \frac{\lambda}{C_\lambda}
        \int_{t>0}\int_{x>0}
        e^{\frac{x}{\lambda} - t}
        P_{X|T_0=t}(x)^2
        dx\,dt
        \,.
    \end{align}
    Recall the bad strip, which we denote by $\mathcal S$,
    \begin{align}
        \mathcal S \triangleq \{(t,x): t\ge 2\epsilon/\lambda,\ |x-\lambda t|\le \epsilon\}.
    \end{align}
    We can continue the lower bound as
    \begin{align}
        \E[R^2]
         & \geq
        \frac{\lambda}{C_\lambda}
        \int_{2\epsilon/\lambda}^\infty \int_{\lambda t - \epsilon}^{\lambda t + \epsilon}
        e^{\frac{x}{\lambda} - t}
        P_{X|T_0=t}(x)^2
        \,dx\,dt
        \\
         & \geq
        \frac{\lambda}{C_\lambda}
        \int_{2\epsilon/\lambda}^\infty \int_{\lambda t - \epsilon}^{\lambda t + \epsilon}
        e^{\frac{\lambda t - \epsilon}{\lambda} - t}
        P_{X|T_0=t}(x)^2
        \,dx\,dt
        \\
         & \geq
        \frac{\lambda}{C_\lambda}
        e^{-\frac{\epsilon}{\lambda}}
        \int_{2\epsilon/\lambda}^\infty \int_{\lambda t - \epsilon}^{\lambda t + \epsilon}
        P_{X|T_0=t}(x)^2
        \,dx\,dt
        \\
         & \geq
        \frac{\lambda}{C_\lambda}
        e^{-\frac{\epsilon}{\lambda}}
        \int_{2\epsilon/\lambda}^\infty \int_{\lambda t - \epsilon}^{\lambda t + \epsilon}
        \inf_{|u|\le \epsilon} p_{X|T_0=t}(u)^2 \,dx\,dt
        \geq
        \infty\,,
    \end{align}
    where to see the last step, first observe that for $\epsilon=\frac12$ there exists $m_\lambda>0$ such that
    \begin{align}
        \inf_{|u|\le \epsilon} p_{X|T_0=t}(u) \ge m_\lambda\,.
    \end{align}
    This uses nothing but the fact that $p_{X|T_0=t}(u)$ is positive everywhere and continuous.

    \paragraph{Step 4: bounding $C_\lambda = \E[e^{|Z|/\lambda}]$.}
    We have
    \begin{align}
        \frac{|Z|}{\lambda}
        \le
        \sum_{k=1}^\infty (1-\lambda)^k |T_k|,
    \end{align}
    hence, by independence,
    \begin{align}
        \E[e^{|Z|/\lambda}]
         & \le
        \prod_{k=1}^\infty \E\!\left[e^{(1-\lambda)^k |T_k|}\right].
    \end{align}
    For a Laplace random variable $L\sim\mathrm{Laplace}(0,1)$ and every $\alpha\in(0,1)$,
    \begin{align}
        \E[e^{\alpha|L|}]
        =
        \int_0^\infty e^{-x}e^{\alpha x}\,dx
        =
        \frac{1}{1-\alpha}.
    \end{align}
    Therefore,
    \begin{align}
        \E[e^{|Z|/\lambda}]
        \le
        \prod_{k=1}^\infty \frac{1}{1-(1-\lambda)^k}.
    \end{align}
    To see that this product is finite, it is enough to show that
    \begin{align}
        \sum_{k=1}^\infty -\log\bigl(1-(1-\lambda)^k\bigr)<\infty.
    \end{align}
    Since $(1-\lambda)^k\to 0$, there exists $K$ such that $(1-\lambda)^k\le 1/2$ for all $k\ge K$.
    For $x\in[0,1/2]$, we have $-\log(1-x)\le 2x$, so for $k\ge K$,
    \begin{align}
        -\log\bigl(1-(1-\lambda)^k\bigr)\le 2(1-\lambda)^k.
    \end{align}
    The right-hand side is summable, so the product is finite and $C_\lambda<\infty$.
\end{proof}

\subsection{Proofs for \cref{sec:log_factor_reduction_to_tail}}

\subsubsection{Proof of \cref{lem:log_factor_fa_quadratic_then_linear}}
\label{proof:log_factor_fa_quadratic_then_linear}
\begin{proof}
    From Taylor's theorem, for $u\le 1/a$ we have
    \begin{align}
        f_a(u)
        =
        e^{-a(u-1)}+a(u-1)-1
        \le
        \frac12 e^a a^2 (u-1)^2.
    \end{align}
    On the event $u>1/a$, we simply use
    \begin{align}
        f_a(u)
        =
        e^{-a(u-1)}+a(u-1)-1
        =
        au
        +
        (e^{-a(u-1)}-a-1)
        \le
        au\,.
    \end{align}
    Therefore, we have
    \begin{align}
        D_{f_a}(P\|Q)
         & =
        \E[f_a(R)]
        \\
         & \le
        e^a a^2 \E[(R-1)^2\mathbf{1}_{\{R\le \frac{1}{a}\}}]
        +
        a \E[R\mathbf{1}_{\{R>\frac{1}{a}\}}]
        \\
         & \le
        e^a
        a^2
        \left(
        1+\E[R_0\mathbf{1}_{\{R_0\le \frac{1}{a}\}}]
        \right)
        +
        a \PP\left(R_0>\frac{1}{a}\right)\,.
        \qedhere
    \end{align}
\end{proof}

\subsubsection{Proof of \cref{lem:log_factor_tail_to_divergence}}
\label{proof:log_factor_tail_to_divergence}
\begin{proof}
    From \cref{eq:log_factor_reduction_to_tail}, we have
    \begin{align}
        D_{f_a}(P\|Q)
         & \le
        e^a
        a^2
        \left(
        1+\E[R_0\mathbf{1}_{\{R_0\le \frac{1}{a}\}}]
        \right)
        +
        a \PP\left(R_0>\frac{1}{a}\right)\,.
    \end{align}
    First, we bound the truncated mean $\E[R_0\mathbf{1}_{\{R_0\le M\}}]$ for a general $M>0$.
    \begin{align}
        \E[R_0\mathbf{1}_{\{R_0\le M\}}]
         & = \int_0^M r p_{R_0}(r) \d r
        \\
         & = \int_0^M \left(\int_{0}^{M} \mathbb{I}_{\{t \le r\}} \d t \right) p_{R_0}(r) \d r
        \\
         & = \int_0^M \left( \int_{0}^{M} \mathbb{I}_{\{t \le r\}}  p_{R_0}(r) \d r \right) \d t
        \\
         & = \int_0^M \left( \int_{t}^{M} p_{R_0}(r) \d r \right) \d t
        \\
         & = \int_0^M \PP\left( t \leq R_0 \leq M \right) \d t
        \\
         & \leq \int_0^M \PP\left( R_0 \geq t \right) \d t
        \\
         & \le
        \int_0^M
        \min\left\{
        1,\frac{C}{t}
        \right\}
        \d t
        \\
         & \le
        C\log(eM)\,,
    \end{align}
    where the second inequality follows from the tail bound in \eqref{eq:log_factor_R0_tail}.
    Plugging $M=1/a$ into the last display gives
    \begin{align}
        D_{f_a}(P\|Q)
         & \le
        e^a a^2
        \left(
        1+C\log\left(\frac{e}{a}\right)
        \right)
        +
        C a^2
        \le
        5C
        a^2 e^a \log\left(\frac{e}{a}\right)\,. \qedhere
    \end{align}
\end{proof}